\documentclass[11pt,oneside]{article}

\usepackage[a4paper,margin=27mm,headheight=15pt]{geometry}
\usepackage[T1]{fontenc}
\usepackage[utf8]{inputenc}
\IfFileExists{libertinus.sty}{\usepackage{libertinus}}{\usepackage{lmodern}}
\usepackage{microtype}
\usepackage{amsmath,amssymb,amsthm,mathtools,mathrsfs}
\usepackage{bm}
\usepackage{booktabs,longtable,array,tabularx}
\usepackage{graphicx}
\usepackage{pgfplots}
\usepackage{xcolor}
\usepackage{enumitem}
\usepackage{fancyhdr}
\usepackage{titlesec}
\usepackage{listings}
\usepackage{xurl}
\usepackage{hyperref}
\usepackage[nameinlink,noabbrev]{cleveref}

\definecolor{linkblue}{RGB}{25,75,135}
\definecolor{shadegray}{RGB}{246,247,249}
\definecolor{rulegray}{RGB}{195,201,208}
\pgfplotsset{compat=1.18}
\hypersetup{
  colorlinks=true,
  linkcolor=linkblue,
  citecolor=linkblue,
  urlcolor=linkblue,
  pdftitle={From Exact Dyadic Formulae to the Small-Argument Asymptotic of the Fabius Function},
  pdfsubject={q-binomial decompression, endpoint-Laplace transfer, and the lower-Lambert phase},
  pdfkeywords={Fabius function, dyadic rational, q-binomial, q-Pochhammer, Thue-Morse, Lambert W, saddle point, asymptotic}
}

\pagestyle{fancy}
\fancyhf{}
\fancyhead[L]{\small Fabius dyadic asymptotics}
\fancyhead[R]{\small\nouppercase{\leftmark}}
\fancyfoot[C]{\thepage}
\renewcommand{\headrulewidth}{0.4pt}

\titleformat{\section}{\Large\bfseries}{\thesection}{0.7em}{}
\titleformat{\subsection}{\large\bfseries}{\thesubsection}{0.7em}{}
\titleformat{\subsubsection}{\normalsize\bfseries}{\thesubsubsection}{0.7em}{}
\setcounter{secnumdepth}{3}
\setcounter{tocdepth}{2}
\setlist[itemize]{leftmargin=2em,itemsep=0.25em,topsep=0.4em}
\setlist[enumerate]{leftmargin=2.3em,itemsep=0.25em,topsep=0.4em}
\setlength{\emergencystretch}{2em}

\newtheorem{theorem}{Theorem}[section]
\newtheorem{proposition}[theorem]{Proposition}
\newtheorem{lemma}[theorem]{Lemma}
\newtheorem{corollary}[theorem]{Corollary}
\newtheorem{conjecture}[theorem]{Conjecture}
\theoremstyle{definition}
\newtheorem{definition}[theorem]{Definition}
\newtheorem{algorithm}[theorem]{Algorithm}
\newtheorem{example}[theorem]{Example}
\theoremstyle{remark}
\newtheorem{remark}[theorem]{Remark}
\newtheorem{warning}[theorem]{Warning}

\newcommand{\R}{\mathbb R}
\newcommand{\C}{\mathbb C}
\newcommand{\Q}{\mathbb Q}
\newcommand{\N}{\mathbb N}
\newcommand{\Z}{\mathbb Z}
\newcommand{\Up}{\operatorname{up}}
\newcommand{\wt}{\operatorname{wt}_2}
\newcommand{\e}{\mathrm e}
\newcommand{\ii}{\mathrm i}
\newcommand{\dd}{\,\mathrm d}
\newcommand{\E}{\mathbb E}
\newcommand{\bigO}{\mathcal O}
\newcommand{\smallo}{o}
\newcommand{\qbinom}[3]{\genfrac{[}{]}{0pt}{}{#1}{#2}_{#3}}
\newcommand{\repo}{\nolinkurl{Analysis/FabiusFunction}}
\newcommand{\M}{\mathfrak m}
\newcommand{\Lam}{\Lambda}
\newcommand{\epsTM}{\varepsilon}
\newcommand{\coeff}[2]{[#1^{#2}]}

\newenvironment{proofidea}{\par\smallskip\noindent\textit{Proof architecture.}\ }{\hfill$\square$\par\smallskip}
\newenvironment{boxedremark}{%
  \par\smallskip\noindent\begin{center}\begin{minipage}{0.94\textwidth}%
  \color{black}\hrule\vspace{0.6em}\small
}{%
  \vspace{0.6em}\hrule\end{minipage}\end{center}\smallskip
}

\lstdefinestyle{pseudo}{
  basicstyle=\ttfamily\small,
  backgroundcolor=\color{shadegray},
  frame=single,
  rulecolor=\color{rulegray},
  columns=fullflexible,
  keepspaces=true,
  showstringspaces=false,
  xleftmargin=0.7em,
  xrightmargin=0.7em,
  aboveskip=0.8em,
  belowskip=0.8em
}

\title{\bfseries From Exact Dyadic Formulae to the Small-Argument\\
Asymptotic of the Fabius Function\\[0.35em]
\large $q$-binomial decompression, endpoint--Laplace transfer, and the lower-Lambert phase}
\author{Technical companion to the formal development\\
\small Vladimir Reshetnikov's \texttt{ProveIt} repository}
\date{Repository snapshot: 25 August 2026}

\begin{document}
\maketitle

\begin{abstract}
The exact arithmetic of the Fabius function and its endpoint asymptotics are usually presented
as two nearly disjoint subjects.  At reciprocal powers of two one has rational recurrences,
composition sums, Thue--Morse power sums, and a finite half-base $q$-binomial formula; near the
origin one has a negative-Laplace product, a lower-Lambert saddle, and a logarithmically
periodic correction.  This report supplies the missing bridge.

The central exact observation is that the apparently formidable $q$-binomial numerator at
$F(2^{-n})$ is not a new asymptotic object.  After choosing a convenient scalar translation,
the inner Thue--Morse block becomes a coefficient of the ratio
$A(-2^kX)/A(-X)$, where
\[
 A(z)=\sum_{n\ge0}a_nz^n
     =\prod_{j\ge1}\frac{\e^{z/2^j}-1}{z/2^j},
 \qquad
 F(2^{-n})=2^{-\binom n2}a_n.
\]
The Gaussian row at $q=1/2$ is exactly the divided-difference functional on the geometric
nodes $-1,-2,\ldots,-2^n$; it annihilates all lower powers and extracts the leading
coefficient.  Consequently the complete finite numerator collapses identically to
$\M_na_n$, where $\M_n=\prod_{j=1}^n(2^j-1)$.  This ``decompression'' proves that the
$q$-Pochhammer, $q$-binomial, Thue--Morse, recurrence, composition, and generating-product
formulae all encode the same coefficient sequence.

Two asymptotic derivations are then developed.  The fully rigorous real-variable route writes
the half moment as $d_n=\E(1-X)^n$, compares it through second order with the negative
Laplace transform $A(-n)=\E\e^{-nX}$, and uses the exact quadratic-plus-periodic
decomposition of $\log A(-s)$.  A complementary coefficient-saddle route applies Cauchy's
formula directly to $a_n=[z^n]A(z)$; it reaches the same lower-Lambert coordinate and provides
an independent explanation of every cancellation.  If $L=\log2$ and $\lambda_n$ is the
large solution of
\[
 \lambda_n-\frac{\log\lambda_n}{L}=n,
 \qquad
 \lambda_n=-\frac1L W_{-1}(-L2^{-n}),
\]
then the resulting sharp dyadic formula is
\[
 \log F(2^{-n})=
 -\frac L2\lambda_n^2+
 \left(1+\frac L2\right)\lambda_n-\frac12\log\lambda_n
 +C_F+\Psi(\lambda_n)+\bigO(n^{-1}),
\]
where $\Psi$ is a smooth, nonconstant, mean-zero, one-periodic function with explicit
Gamma--zeta Fourier coefficients.  Expanding the Lambert phase yields
\[
 \log F(2^{-n})=
 -\frac L2n^2-n\log n+\left(1+\frac L2\right)n
 -\frac{\log^2n}{2L}+C_F-
 \frac{\log^2n}{2L^2n}+\Psi(\lambda_n)+\bigO(n^{-1}).
\]
The exact phase is essential: replacing it by $\log_2n$ leaves a generally larger
$(\log n)/n$ term.  Corollaries give equally sharp asymptotics for the recurrence sequence,
the half moments, the composition sum, and the literal $q$-binomial numerator.  Numerical
checks against exact rational values confirm the constant, phase, and first saddle
correction.
\end{abstract}

\begingroup\footnotesize
\tableofcontents
\endgroup
\newpage

\section{The missing bridge}
\label{sec:introduction}

The Fabius function admits two descriptions that look, at first sight, almost antagonistic.
Its values at dyadic rationals are algebraic and finite: every value is rational; there are
terminating recurrences; the Thue--Morse sequence provides finite signed power sums; and a
half-base Gaussian-binomial row packages those sums into a compact formula.  Its endpoint
asymptotic is analytic and infinite: one studies a Laplace product, separates a periodic
renormalization of a dyadic harmonic sum, and evaluates a saddle determined by the lower real
branch of Lambert's $W$-function.

The two subjects are not merely compatible.  They are two coordinate systems on the same
object.  The exact formulae contain the asymptotic information in a highly compressed form,
and the purpose of this report is to make the decompression explicit.

\subsection{What is proved}

There are four levels of connection.

\begin{enumerate}
\item The reciprocal-power value $F(2^{-n})$ is exactly a half moment divided by a factorial
and a triangular power of two.  This already gives the quadratic logarithmic decay.

\item After factorial normalization, all reciprocal-power values are coefficients of one
entire function $A$.  The recurrence, composition sum, and infinite product are equivalent
presentations of $A$.

\item The finite $q$-binomial/Thue--Morse formula can be reduced identically to the same
coefficient $a_n=[z^n]A(z)$.  No limiting argument is involved.  The $q$-binomial row is a
geometric divided difference; its role is to undo the blockwise Thue--Morse cancellation.

\item The product for $A(-s)$ has an exact quadratic-plus-periodic logarithm.  Transferring
from $A(-n)$ to the endpoint moment, or applying a coefficient saddle directly to $A$, yields
the corrected lower-Lambert asymptotic.
\end{enumerate}

The first, second, third, and the full real-variable fourth level are represented in the Lean
formalization under \repo.  The coefficient-saddle calculation in \Cref{sec:coefficient-saddle}
is included as an independent conventional-mathematics derivation; it is useful because it
starts from the coefficient formula itself and exposes the same cancellations with very
little probability language.

\subsection{Headline notation and result}

Throughout,
\begin{equation}
 L=\log2.                                                     \label{eq:L}
\end{equation}
Let $P$ be the periodic correction in the exact negative-Laplace decomposition, let
\begin{equation}
 \overline P=\int_0^1P(u)\dd u,
 \qquad \Psi(u)=P(u)-\overline P,                            \label{eq:Psi-intro}
\end{equation}
and put
\begin{equation}
 A_0=\frac{6\gamma^2+12\gamma_1-\pi^2}{12},
 \qquad
 C_F=\frac{A_0}{L}-\frac{7L}{12}-\frac12\log\pi.            \label{eq:CF-intro}
\end{equation}
Here $\gamma$ is Euler's constant and $\gamma_1$ is the first Stieltjes constant.
For every sufficiently large integer $n$, let $\lambda_n$ be the large positive solution of
\begin{equation}
 \lambda_n2^{-\lambda_n}=2^{-n},
 \qquad\text{equivalently}\qquad
 \lambda_n-\frac{\log\lambda_n}{L}=n.                       \label{eq:lambda-n-intro}
\end{equation}

\begin{theorem}[Sharp reciprocal-power asymptotic]
\label{thm:headline}
As $n\to\infty$,
\begin{equation}
 \boxed{
 \log F(2^{-n})=
 -\frac L2\lambda_n^2+
 \left(1+\frac L2\right)\lambda_n-\frac12\log\lambda_n
 +C_F+\Psi(\lambda_n)+\bigO(n^{-1}).}                       \label{eq:headline-compact}
\end{equation}
Equivalently,
\begin{align}
 \log F(2^{-n})={}&-\frac L2n^2-n\log n+
 \left(1+\frac L2\right)n-\frac{\log^2n}{2L}+C_F           \notag\\
 &-\frac{\log^2n}{2L^2n}+\Psi(\lambda_n)+\bigO(n^{-1}).   \label{eq:headline-expanded}
\end{align}
The function $\Psi$ is smooth, mean-zero, one-periodic, and nonconstant.  Its Fourier
coefficients are explicit.
\end{theorem}

The rest of the report derives \Cref{thm:headline} from the exact arithmetic formulae rather
than taking it as an externally supplied endpoint theorem.

\begin{boxedremark}
The most important conceptual point is that the periodic term is already present in the exact
dyadic values.  It does not arise because one interpolates between dyadic points.  The
fractional parts of the exact saddle coordinates $\lambda_n$ drift through the logarithmic
period, so the single rational sequence $F(2^{-n})$ samples the entire periodic correction.
\end{boxedremark}

\section{One coefficient sequence behind the exact formulae}
\label{sec:exact-sequence}

\subsection{The half moments and the endpoint identity}

Let $F$ denote the bounded Fabius distribution function and let $\Up$ be Rvachev's even
compactly supported density in the normalization
\begin{equation}
 \Up(t)=F(1-|t|),\qquad |t|\le1.                             \label{eq:up-fold}
\end{equation}
Define the half moments
\begin{equation}
 d_0=1,
 \qquad
 d_n=n\int_0^1t^{n-1}\Up(t)\dd t\quad(n\ge1).               \label{eq:dn}
\end{equation}
They are rational and positive.  Their first few values are
\begin{equation}
 d_0=1,\qquad d_1=\frac12,\qquad d_2=\frac5{18},\qquad
 d_3=\frac16,\qquad d_4=\frac{143}{1350}.                  \label{eq:dn-first}
\end{equation}

\begin{theorem}[Exact endpoint--moment bridge]
\label{thm:endpoint-moment}
For every $n\ge0$,
\begin{equation}
 d_n=n!\,2^{\binom n2}F(2^{-n}),
 \qquad
 F(2^{-n})=\frac{d_n}{n!\,2^{\binom n2}}.                  \label{eq:endpoint-moment}
\end{equation}
\end{theorem}

\begin{proof}
The primitive-chain proof is especially revealing.  For $r\ge0$ define
\begin{equation}
 f_r(t)=2^{\binom r2}\Up\bigl((t+1)2^{-r}-1\bigr),
 \qquad -1\le t\le1.                                       \label{eq:primitive-chain}
\end{equation}
The functional--differential equation of $\Up$ implies
\begin{equation}
 f_{r+1}'(t)=f_r(t),\qquad f_r(-1)=0.                        \label{eq:primitive-chain-derivative}
\end{equation}
Put $J_n=\int_{-1}^0t^{n-1}\Up(t)\dd t$.  Repeated integration by parts, using
\eqref{eq:primitive-chain-derivative}, gives
\begin{equation}
 J_n=(-1)^{n-1}(n-1)!f_n(0).                                \label{eq:J-chain}
\end{equation}
Evenness of $\Up$ turns the left side into
$(-1)^{n-1}\int_0^1t^{n-1}\Up(t)\dd t$, while
\begin{equation*}
 f_n(0)=2^{\binom n2}\Up(2^{-n}-1)
       =2^{\binom n2}F(2^{-n}).
\end{equation*}
Multiplication by $n$ proves \eqref{eq:endpoint-moment}.  The case $n=0$ is
$F(1)=d_0=1$.
\end{proof}

There is also a probability interpretation that will be decisive later.  Let $X$ be a random
variable with distribution function $F$.  Integrating by parts against the distribution
measure gives
\begin{equation}
 d_n=\E(1-X)^n.                                              \label{eq:dn-endpoint-probability}
\end{equation}
Thus $d_n$ is simultaneously an exact rational endpoint value and an endpoint moment of the
probability law.

\subsection{Factorial normalization and the recurrence}

Set
\begin{equation}
 a_n=\frac{d_n}{n!},
 \qquad
 A(z)=\sum_{n=0}^{\infty}a_nz^n.                            \label{eq:an-A}
\end{equation}
Then \Cref{thm:endpoint-moment} becomes the particularly economical identity
\begin{equation}
 \boxed{F(2^{-n})=2^{-\binom n2}a_n.}                       \label{eq:F-an}
\end{equation}
The half-moment recurrence is equivalent to
\begin{equation}
 (2^n-1)a_n=\sum_{k=0}^{n-1}\frac{a_k}{(n-k+1)!},
 \qquad n\ge1,
 \qquad a_0=1.                                              \label{eq:an-recurrence}
\end{equation}
Equivalently, directly in terms of endpoint values,
\begin{equation}
 F(2^{-n})=
 \frac{2^{-\binom n2}}{2^n-1}
 \sum_{k=0}^{n-1}\frac{2^{\binom k2}}{(n-k+1)!}F(2^{-k}).  \label{eq:F-recurrence}
\end{equation}

The recurrence is not merely an evaluation algorithm.  It is already a coefficient equation.
Let
\begin{equation}
 E(z)=\frac{\e^z-1}{z},\qquad E(0)=1.                       \label{eq:E}
\end{equation}
Multiplying \eqref{eq:an-recurrence} by $z^n$ and summing gives
\begin{equation}
 \boxed{A(2z)=E(z)A(z).}                                    \label{eq:A-functional}
\end{equation}
Conversely, coefficient comparison in \eqref{eq:A-functional} recovers
\eqref{eq:an-recurrence}.  Thus the exact recurrence and the analytic refinement equation are
the same statement.

\subsection{Composition sum}

Repeated substitution in \eqref{eq:an-recurrence} gives a finite closed form.  A composition
$\rho\vDash n$ is an ordered tuple $\rho=(r_1,\ldots,r_m)$ of positive integers with sum
$n$; write $s_j=r_1+\cdots+r_j$ for its prefix sums.

\begin{proposition}[Exact composition formula]
\label{prop:composition}
For $n\ge1$,
\begin{equation}
 a_n=
 \sum_{\rho\vDash n}
 \prod_{j=1}^{\ell(\rho)}
 \frac{1}{(2^{s_j}-1)(r_j+1)!}.                             \label{eq:composition-an}
\end{equation}
Consequently
\begin{equation}
 F(2^{-n})=2^{-\binom n2}
 \sum_{\rho\vDash n}
 \prod_{j=1}^{\ell(\rho)}
 \frac{1}{(2^{s_j}-1)(r_j+1)!}.                             \label{eq:composition-F}
\end{equation}
\end{proposition}

\begin{proof}
Every use of \eqref{eq:an-recurrence} replaces a current index $s_j$ by a smaller prefix
$s_{j-1}$ and contributes
$((2^{s_j}-1)(s_j-s_{j-1}+1)!)^{-1}$.  A complete descent from $n$ to $0$ is therefore the
same as a composition of $n$, with increments $r_j=s_j-s_{j-1}$.  Summing over all descents
gives \eqref{eq:composition-an}.
\end{proof}

This formula looks combinatorial, but asymptotically it is not a sum over independently
comparable contributions.  The number of compositions is $2^{n-1}$, and the dominant family
moves with $n$.  The correct way to analyze the sum is to reassemble it into $A$.

\subsection{The entire product}

Iterating \eqref{eq:A-functional} toward the origin gives
\begin{equation}
 \boxed{
 A(z)=\prod_{j=1}^{\infty}E(z/2^j)
 =\prod_{j=1}^{\infty}\frac{\e^{z/2^j}-1}{z/2^j}.}         \label{eq:A-product}
\end{equation}
The product converges locally uniformly and defines an entire function.  It is the moment
generating function of $X$:
\begin{equation}
 A(z)=\E\e^{zX}.                                            \label{eq:A-mgf}
\end{equation}
Since $E(z)=\e^zE(-z)$ and $\sum_{j\ge1}2^{-j}=1$, the product also satisfies the exact
reflection identity
\begin{equation}
 A(z)=\e^zA(-z).                                             \label{eq:A-reflection}
\end{equation}

The equivalence graph is now
\begin{equation*}
 \boxed{\begin{gathered}
 \text{endpoint values}\longleftrightarrow d_n\longleftrightarrow a_n
 \longleftrightarrow\text{recurrence}\\
 \longleftrightarrow\text{composition sum}\longleftrightarrow A(z)
 \longleftrightarrow\text{entire product}
 \end{gathered}}
\end{equation*}
The next section inserts the $q$-binomial and Thue--Morse formula into this same graph.

\section{Decompressing the half-base \texorpdfstring{$q$}{q}-binomial formula}
\label{sec:q-decompression}

The $q$-binomial formula is the most conspicuous exact representation and the least transparent
asymptotic one.  Its finite sum contains two distinct annihilation mechanisms:

\begin{itemize}
\item inside each dyadic block, Thue--Morse signs annihilate low polynomial degrees;
\item across the block sizes, a Gaussian-binomial divided difference annihilates the remaining
low powers of the geometric scale.
\end{itemize}

Once both cancellations are written as coefficient extraction, the entire numerator reduces
to $\M_na_n$.

\subsection{The geometric interpolation functional}

For $q=1/2$, write
\begin{equation}
 (a;1/2)_n=\prod_{j=0}^{n-1}(1-a2^{-j}),                    \label{eq:qpoch}
\end{equation}
and let $\qbinom nk{1/2}$ denote the Gaussian binomial coefficient
\cite{andrews1976,gasper2004,kac2002,dlmf17}.  The finite
$q$-binomial theorem is
\begin{equation}
 (z;1/2)_n=\sum_{k=0}^n
 (-1)^k2^{-\binom k2}\qbinom nk{1/2}z^k.                   \label{eq:q-binomial-theorem}
\end{equation}
Define
\begin{equation}
 x_k=-2^k,
 \qquad
 w_{n,k}=(-1)^k2^{-\binom k2}\qbinom nk{1/2},
 \qquad
 \M_n=\prod_{j=1}^n(2^j-1).                                \label{eq:q-weights}
\end{equation}

\begin{lemma}[Leading-coefficient extraction]
\label{lem:q-leading}
If $R$ is a polynomial of degree at most $n$, then
\begin{equation}
 \sum_{k=0}^nw_{n,k}R(-2^k)=\M_n[z^n]R(z).                 \label{eq:q-leading}
\end{equation}
Equivalently,
\begin{equation}
 \sum_{k=0}^nw_{n,k}x_k^d=
 \begin{cases}
 0,&0\le d<n,\\
 \M_n,&d=n.
 \end{cases}                                                \label{eq:q-monomials}
\end{equation}
\end{lemma}

\begin{proof}
The normalized weights are the barycentric weights for the nodes
$x_0,\ldots,x_n$:
\begin{equation}
 \frac{w_{n,k}}{\M_n}=\frac1{\prod_{j\ne k}(x_k-x_j)}.      \label{eq:barycentric}
\end{equation}
The coefficient of $z^n$ in the Lagrange interpolation formula for $R$ is therefore the left
side of \eqref{eq:q-leading} divided by $\M_n$.  Formula \eqref{eq:barycentric} can also be
checked directly from \eqref{eq:q-binomial-theorem}.
\end{proof}

The Pochhammer normalization is exactly the same Mersenne product:
\begin{equation}
 (1/2;1/2)_n=2^{-\binom{n+1}{2}}\M_n.                      \label{eq:q-mersenne}
\end{equation}
This identity will make the last power-of-two cancellation automatic.

\subsection{The finite Thue--Morse blocks}

Let
\begin{equation}
 \epsTM_r=(-1)^{\wt(r)}.                                    \label{eq:tm-sign}
\end{equation}
For $k,N\ge0$ and a scalar translation $Q$, define
\begin{equation}
 S_{k,N}(Q)=\sum_{r=0}^{2^k-1}\epsTM_r(r-2^k+Q)^N.          \label{eq:Skn}
\end{equation}
The complete-block exponential generating function is
\begin{equation}
 C_k(X)=\sum_{r<2^k}\epsTM_r\e^{(r-2^k)X}.                 \label{eq:Ck}
\end{equation}
The elementary binary product
\begin{equation}
 \sum_{r<2^k}\epsTM_rz^r=\prod_{j=0}^{k-1}(1-z^{2^j})     \label{eq:tm-product}
\end{equation}
implies
\begin{equation}
 \frac{C_k(X)}{X^k}
 =(-1)^k2^{\binom k2}\e^{-X}
   \prod_{j=0}^{k-1}E(-2^jX).                              \label{eq:Ck-factor}
\end{equation}
In particular, $C_k$ has a zero of order $k$ at the origin.  Equivalently,
$S_{k,N}(Q)=0$ for $N<k$, independently of $Q$.

Iterating \eqref{eq:A-functional} on the negative dyadic orbit gives
\begin{equation}
 \frac{A(-2^kX)}{A(-X)}=\prod_{j=0}^{k-1}E(-2^jX).          \label{eq:A-ratio}
\end{equation}
Translation invariance permits the especially convenient choice $Q=1$.  Since
\begin{equation}
 \sum_{N\ge0}S_{k,N}(1)\frac{X^N}{N!}=\e^XC_k(X),          \label{eq:S-egf}
\end{equation}
\eqref{eq:Ck-factor} and \eqref{eq:A-ratio} yield the exact coefficient identity
\begin{equation}
 \boxed{
 \frac{S_{k,n+k}(1)}{(n+k)!}
 =(-1)^k2^{\binom k2}
 [X^n]\frac{A(-2^kX)}{A(-X)}.}                             \label{eq:S-ratio-coeff}
\end{equation}
This is the first decompression: an inner Thue--Morse power sum is a coefficient of the same
entire function that generates the half moments.

\subsection{The literal numerator collapses}

At $m=1$, the global $q$-binomial formula reads
\begin{equation}
 F(2^{-n})=\frac{\mathcal N_n}
 {2^{n^2}(1/2;1/2)_n},                                     \label{eq:q-F-N}
\end{equation}
where, using the shift $Q=1$,
\begin{equation}
 \mathcal N_n=
 \sum_{k=0}^n
 \frac{\qbinom nk{1/2}}{4^{\binom k2}(n+k)!}
 S_{k,n+k}(1).                                              \label{eq:q-numerator}
\end{equation}

\begin{theorem}[$q$-binomial decompression]
\label{thm:q-decompression}
For every $n\ge0$,
\begin{equation}
 \boxed{\mathcal N_n=\M_na_n.}                             \label{eq:N-m-a}
\end{equation}
Consequently the finite $q$-binomial formula is identically the endpoint formula
$F(2^{-n})=2^{-\binom n2}a_n$.
\end{theorem}

\begin{proof}
For fixed $n$, introduce the polynomial
\begin{equation}
 R_n(z)=[X^n]\frac{A(zX)}{A(-X)}.                           \label{eq:Rn}
\end{equation}
It has degree at most $n$.  Since $A(-X)^{-1}=1+\bigO(X)$ and
$A(zX)=\sum_{j\ge0}a_jz^jX^j$, its leading coefficient is
\begin{equation}
 [z^n]R_n(z)=a_n.                                           \label{eq:Rn-leading}
\end{equation}
Substitution of \eqref{eq:S-ratio-coeff} into \eqref{eq:q-numerator} gives
\begin{align}
 \mathcal N_n
 &=\sum_{k=0}^n
 \qbinom nk{1/2}2^{-2\binom k2}
 (-1)^k2^{\binom k2}R_n(-2^k)                              \notag\\
 &=\sum_{k=0}^n w_{n,k}R_n(-2^k).                          \label{eq:N-interpolation}
\end{align}
By \Cref{lem:q-leading}, this is $\M_n[z^n]R_n(z)=\M_na_n$, proving
\eqref{eq:N-m-a}.  Finally, \eqref{eq:q-mersenne} gives
\begin{equation*}
 \frac{\M_na_n}{2^{n^2}(1/2;1/2)_n}
 =\frac{\M_na_n}{2^{n^2-\binom{n+1}{2}}\M_n}
 =2^{-\binom n2}a_n.
\end{equation*}
\end{proof}

\begin{boxedremark}
The two finite cancellations have now been separated completely.  The Thue--Morse block
removes $k$ powers of $X$ and converts a power sum into $R_n(-2^k)$.  The half-base
$q$-binomial row then removes all powers of the scale node below degree $n$ and returns the
coefficient $a_n$.  The $q$-Pochhammer factor does not introduce a new special-function
asymptotic; it is the normalization of this divided difference.
\end{boxedremark}

\subsection{General dyadic numerators}

For completeness, the exact global formula is
\begin{align}
 \mathcal F(m/2^n)
 &=\frac{1}{2^{n^2}(1/2;1/2)_n}
 \sum_{k=0}^{n}
 \frac{\qbinom nk{1/2}}{4^{\binom k2}(n+k)!}                \notag\\
 &\qquad\times
 \sum_{r=0}^{m2^k-1}\epsTM_r(r-m2^k+Q)^{n+k}.              \label{eq:q-global}
\end{align}
It is independent of $Q$.  Splitting $r=h2^k+s$ factors the long block into the complete
block $C_k$ and a finite prefix polynomial.  The same interpolation functional then extracts
one coefficient of a product of that prefix polynomial with $A$.  Thus even the general
formula contains no asymptotically unrelated object.  What changes when $m$ varies with $n$
is the coefficient multiplier, not the underlying dyadic product.

\section{The first asymptotic already visible in the exact values}
\label{sec:coarse}

Before resolving the periodic fluctuation, the endpoint identity alone gives the principal
quadratic law.  Taking logarithms in \eqref{eq:endpoint-moment},
\begin{equation}
 \log F(2^{-n})=\log d_n-\log(n!)-\binom n2L.               \label{eq:log-F-exact}
\end{equation}
Because $0\le\Up\le1$,
\begin{equation}
 d_n\le n\int_0^1t^{n-1}\dd t=1.                           \label{eq:dn-upper}
\end{equation}
A fixed positive portion of the mass on $[0,1/2]$ gives, for $n\ge1$,
\begin{equation}
 2^{-(n+1)}\le d_n.                                         \label{eq:dn-lower}
\end{equation}
Combining these bounds with elementary factorial estimates yields the explicit inequality
\begin{equation}
 \left|\log F(2^{-n})+\frac L2n^2\right|
 \le3n\log(n+1),\qquad n\ge1.                              \label{eq:coarse-bound}
\end{equation}
Therefore
\begin{equation}
 \lim_{n\to\infty}\frac{\log F(2^{-n})}{n^2}=-\frac L2.    \label{eq:coarse-limit}
\end{equation}

This is already a derivation of the leading endpoint asymptotic from exact arithmetic.  It
also explains why neither the factorial nor the half moment can change the coefficient of
$n^2$: $\log(n!)=\bigO(n\log n)$ and $\log d_n=\bigO(n)$ under the crude bounds.  To recover
the $-n\log n$ term, the $\log^2n$ term, the constant, and the periodic fluctuation, however,
one must understand $d_n$ much more sharply.

\section{The exact negative-Laplace decomposition}
\label{sec:laplace-product}

The sharp asymptotic enters through the same function $A$.  For $s>0$ put
\begin{equation}
 B(s)=A(-s)=\E\e^{-sX},
 \qquad
 \Lam(s)=\log B(s).                                         \label{eq:B-Lambda}
\end{equation}
From \eqref{eq:A-product},
\begin{equation}
 B(s)=\prod_{j=1}^{\infty}
 \frac{1-\e^{-s/2^j}}{s/2^j},                               \label{eq:B-product}
\end{equation}
and hence
\begin{equation}
 \Lam(s)=\sum_{j=1}^{\infty}\kappa(s/2^j),
 \qquad
 \kappa(u)=\log\frac{1-\e^{-u}}u.                          \label{eq:Lambda-harmonic}
\end{equation}
This is a dyadic harmonic sum.  Its dilation equation is exact:
\begin{equation}
 \Lam(2s)=\Lam(s)+\log(1-\e^{-s})-\log s.                  \label{eq:Lambda-dilation}
\end{equation}

\subsection{Periodization of the dilation equation}

Define the forward logarithmic tail
\begin{equation}
 T(s)=\sum_{m=0}^{\infty}\log(1-\e^{-s2^m})<0,
 \qquad
 \mathcal E(s)=-T(s)>0.                                    \label{eq:tail}
\end{equation}
Since $2^m\ge m+1$,
\begin{equation}
 0\le\mathcal E(s)\le
 \frac{\e^{-s}}{(1-\e^{-s})^2},                            \label{eq:tail-bound-report}
\end{equation}
and in particular $\mathcal E(s)\le4\e^{-s}$ for $s\ge L$.
For $u\in\R$, set
\begin{equation}
 P(u)=\Lam(2^u)+\frac L2(u^2-u)+T(2^u).                    \label{eq:P-report}
\end{equation}

\begin{theorem}[Exact quadratic-plus-periodic decomposition]
\label{thm:periodic-decomposition}
The function $P$ is smooth and one-periodic.  For every $s>0$,
\begin{equation}
 \boxed{
 \Lam(s)=-\frac{(\log s)^2}{2L}+\frac12\log s
 +P(\log_2s)+\mathcal E(s).}                               \label{eq:Lambda-periodic}
\end{equation}
All derivatives of $P$ are bounded.
\end{theorem}

\begin{proof}
The forward tail satisfies
\begin{equation*}
 T(2s)=T(s)-\log(1-\e^{-s}).
\end{equation*}
Under $u\mapsto u+1$, the change in the quadratic term of \eqref{eq:P-report} is
\begin{equation*}
 \frac L2\bigl((u+1)^2-(u+1)-(u^2-u)\bigr)=Lu=\log s.
\end{equation*}
This cancels the changes in \eqref{eq:Lambda-dilation} and in $T$, proving
$P(u+1)=P(u)$.  Rearrangement gives \eqref{eq:Lambda-periodic}.  Local uniform convergence of
the defining series permits differentiation; periodicity then makes every derivative bounded.
\end{proof}

The drift $-(\log s)^2/(2L)$ is the analytic counterpart of the triangular power
$2^{-\binom n2}$ in the exact endpoint formula.  The periodic remainder is the residual
freedom in solving the dilation equation on logarithmic scale.

\subsection{Mean, Fourier modes, and the asymptotic constant}

Let $\overline P$ and $\Psi$ be as in \eqref{eq:Psi-intro}.  The Mellin transform of the
kernel gives
\begin{equation}
 \overline P=\frac{A_0}{L}-\frac L{12},                     \label{eq:P-mean}
\end{equation}
where $A_0$ is defined in \eqref{eq:CF-intro}.  With
\begin{equation}
 \chi_k=\frac{2\pi\ii k}{L},                                \label{eq:chi}
\end{equation}
the Fourier coefficients of the centered function are
\begin{equation}
 \widehat\Psi(0)=0,
 \qquad
 \boxed{
 \widehat\Psi(k)=-\frac{\Gamma(-\chi_k)\zeta(1-\chi_k)}L}
 \quad(k\ne0).                                             \label{eq:Psi-Fourier-report}
\end{equation}
The Gamma factor has no zeros, and $\zeta$ has no zeros on $\Re z=1$ away from its pole at
$1$, so every nonzero Fourier coefficient is nonzero.  Thus $\Psi$ is genuinely nonconstant.
Stirling's vertical estimate for $\Gamma$ makes the coefficients exponentially small; in
fact the oscillation has amplitude only a few parts in $10^6$, but it cannot be deleted at an
error scale tending to zero.

The Gaussian normalization later contributes $-\tfrac12\log(2\pi)$.  Therefore
\begin{equation}
 \overline P-\frac12\log(2\pi)
 =\frac{A_0}{L}-\frac{7L}{12}-\frac12\log\pi
 =C_F.                                                       \label{eq:CF-from-P}
\end{equation}
Numerically,
\begin{equation}
 \overline P=-1.10904536543418668\ldots,
 \qquad
 C_F=-2.02798389863885942\ldots.                            \label{eq:constants-numeric}
\end{equation}

\section{Sharp route I: endpoint moments versus Laplace moments}
\label{sec:endpoint-laplace}

This route begins with the exact rational object $d_n=\E(1-X)^n$ and converts it into the
negative-Laplace product at the natural scale $s=n$.  It is the route closest to the current
Lean bridge modules.

\subsection{Second-order kernel comparison}

For $0\le x\le1$ and $n\ge1$, the logarithmic expansion
\begin{equation*}
 n\log(1-x)=-nx-\frac n2x^2+\bigO(nx^3)
\end{equation*}
is useful only after the large-$x$ region is controlled uniformly.  A convenient global
form is
\begin{equation}
 \left|(1-x)^n-\e^{-nx}\left(1-\frac n2x^2\right)\right|
 \le16\e^{-nx}(nx^3+n^2x^4).                               \label{eq:kernel-comparison}
\end{equation}
Integrating against the law of $X$ gives
\begin{align}
 d_n={}&B(n)-\frac n2\E\bigl(X^2\e^{-nX}\bigr)+R_n,        \label{eq:dn-B-second}\\
 |R_n|\le{}&16\left(
 n\E(X^3\e^{-nX})+n^2\E(X^4\e^{-nX})\right).              \label{eq:Rn-bound}
\end{align}

Let $q(s)=\Lam(s)=\log B(s)$.  The normalized tilted second moment is
\begin{equation}
 \frac{\E(X^2\e^{-sX})}{B(s)}=q''(s)+q'(s)^2.               \label{eq:tilted-second}
\end{equation}
The corresponding third- and fourth-moment identities are polynomials in
$q',q'',q^{(3)},q^{(4)}$.  Differentiation of \eqref{eq:Lambda-periodic}, together with the
exponentially small tail, shows that these normalized moments have precisely the inverse
powers needed to make \eqref{eq:Rn-bound} a relative $\bigO(n^{-1})$ error after the second
term is retained.

\begin{proposition}[Endpoint--Laplace logarithmic transfer]
\label{prop:endpoint-laplace-transfer}
As $n\to\infty$,
\begin{equation}
 \boxed{
 \log d_n=q(n)-\frac n2\bigl(q''(n)+q'(n)^2\bigr)
 +\bigO(n^{-1}).}                                          \label{eq:log-d-transfer}
\end{equation}
\end{proposition}

\begin{proofidea}
Divide \eqref{eq:dn-B-second} by $B(n)$.  The second-order displacement is
\begin{equation*}
 \alpha_n=\frac n2\bigl(q''(n)+q'(n)^2\bigr)
          =\bigO(\log^2n/n),
\end{equation*}
and the normalized remainder is $\bigO(n^{-1})$.  Thus
$d_n/B(n)=1-\alpha_n+\bigO(n^{-1})$.  The local estimate
$\log(1-y)=-y+\bigO(y^2)$ applies; since
$\alpha_n^2=\bigO(\log^4n/n^2)=\bigO(n^{-1})$, it gives
\eqref{eq:log-d-transfer}.
\end{proofidea}

\subsection{Differentiating the exact periodic decomposition}

Write
\begin{equation}
 \ell=\log s,
 \qquad
 u=\log_2s=\frac\ell L.                                    \label{eq:ell-u}
\end{equation}
Ignoring an exponentially small differentiated tail, \eqref{eq:Lambda-periodic} gives
\begin{equation}
 q'(s)=\frac{-\ell/L+1/2+\Psi'(u)/L}{s}+\bigO(\e^{-s}),     \label{eq:qprime}
\end{equation}
and
\begin{equation}
 q''(s)=\frac{\ell/L-1/2-1/L-\Psi'(u)/L+\Psi''(u)/L^2}{s^2}
 +\bigO(\e^{-s}).                                           \label{eq:qsecond}
\end{equation}
Squaring \eqref{eq:qprime} and adding \eqref{eq:qsecond}, the terms linear in $\ell$ that do
not involve $\Psi'$ cancel.  Since $\Psi'$ and $\Psi''$ are bounded,
\begin{equation}
 \frac s2\bigl(q''(s)+q'(s)^2\bigr)
 =\frac{\ell^2}{2L^2s}-
   \frac{\ell\,\Psi'(u)}{L^2s}+\bigO(s^{-1}).              \label{eq:cumulant-expansion}
\end{equation}
This cancellation is the analytic reason that the first nonperiodic endpoint correction is
$-\log^2n/(2L^2n)$ rather than a mixture of $\log^2n/n$ and $\log n/n$ terms.

Substituting \eqref{eq:Lambda-periodic} and \eqref{eq:cumulant-expansion} into
\eqref{eq:log-d-transfer} yields the first sharp form.

\begin{theorem}[Half moments at the logarithmic phase]
\label{thm:dn-logphase}
Let $u_n=\log_2n$.  Then
\begin{align}
 \log d_n={}&-\frac{\log^2n}{2L}+\frac12\log n
 +\overline P+\Psi(u_n)                                    \notag\\
 &-\frac{\log^2n}{2L^2n}
 +\frac{\log n}{L^2n}\Psi'(u_n)+\bigO(n^{-1}).             \label{eq:dn-logphase}
\end{align}
\end{theorem}

Combining \eqref{eq:dn-logphase} with Stirling's formula and
\eqref{eq:log-F-exact} gives
\begin{align}
 \log F(2^{-n})={}&-\frac L2n^2-n\log n+
 \left(1+\frac L2\right)n-\frac{\log^2n}{2L}+C_F           \notag\\
 &-\frac{\log^2n}{2L^2n}+\Psi(u_n)
 +\frac{\log n}{L^2n}\Psi'(u_n)+\bigO(n^{-1}).             \label{eq:F-periodic-derivative}
\end{align}
Formula \eqref{eq:F-periodic-derivative} has already recovered the constant and the periodic
function from the exact endpoint values.  The derivative term is not an extra oscillation; it
is the first-order displacement from the naive phase $u_n$ to the exact Lambert phase.

\section{The lower-Lambert phase and the derivative cancellation}
\label{sec:lambert-phase}

\subsection{The exact phase}

The natural saddle coordinate at $x=2^{-n}$ is
\begin{equation}
 \lambda_n=-\frac1L W_{-1}(-L2^{-n}),                       \label{eq:lambda-W}
\end{equation}
where $W_{-1}$ is the lower real branch of Lambert's function
\cite{corless1996}.  It is the large solution of \eqref{eq:lambda-n-intro}.  Its fixed-point
form is
\begin{equation}
 \lambda_n=n+\log_2\lambda_n.                               \label{eq:lambda-fixed}
\end{equation}
The standard lower-Lambert expansion gives
\begin{align}
 \lambda_n={}&n+\frac{\log n}{L}+\frac{\log n}{L^2n}
 +\frac{\log n-\tfrac12\log^2n}{L^3n^2}                    \notag\\
 &+\frac{\log n-\tfrac32\log^2n+\tfrac13\log^3n}{L^4n^3}
 +\bigO\!\left(\frac{\log^4n}{n^4}\right).               \label{eq:lambda-expansion}
\end{align}
For the periodic phase transfer only the first displacement is needed.  Define
\begin{equation}
 \delta_n=\log_2\lambda_n-\log_2n.                          \label{eq:delta}
\end{equation}
Then
\begin{equation}
 \delta_n=\frac{\log n}{L^2n}
 +\bigO\!\left(\frac{\log^2n}{n^2}\right).                \label{eq:delta-expansion}
\end{equation}

\subsection{Why \texorpdfstring{$\Psi$}{Psi} is evaluated at \texorpdfstring{$\lambda_n$}{lambda n}}

Since $n$ is an integer, periodicity and \eqref{eq:lambda-fixed} imply the exact identity
\begin{equation}
 \Psi(\lambda_n)=\Psi(\lambda_n-n)
 =\Psi(\log_2\lambda_n).                                   \label{eq:Psi-lambda-exact}
\end{equation}
Taylor expansion at $u_n=\log_2n$, with uniformly bounded second derivative, gives
\begin{align}
 \Psi(\lambda_n)
 &=\Psi(u_n+\delta_n)                                      \notag\\
 &=\Psi(u_n)+\frac{\log n}{L^2n}\Psi'(u_n)
   +\bigO(n^{-1}).                                          \label{eq:Psi-phase-transfer}
\end{align}
Indeed, the error in \eqref{eq:delta-expansion} contributes
$\bigO(\log^2n/n^2)=\bigO(n^{-1})$, and the quadratic Taylor remainder contributes the same
or less.  Substitution of \eqref{eq:Psi-phase-transfer} into
\eqref{eq:F-periodic-derivative} proves \eqref{eq:headline-expanded}.

This calculation explains exactly why the phase may not be replaced by $\log_2n$ at the
claimed precision.  One would omit the term
\begin{equation}
 \frac{\log n}{L^2n}\Psi'(\log_2n),                         \label{eq:missing-phase-term}
\end{equation}
which is generally of order $(\log n)/n$, not $1/n$.  The tiny amplitude of $\Psi$ affects
numerical visibility, but not the asymptotic order.

\begin{warning}[Phase precision is amplified]
The leading exponent contains $-L\lambda^2/2$.  An error $\Delta\lambda$ in the phase changes
the exponent by approximately $-L\lambda\Delta\lambda$.  Therefore deriving an elementary
expansion with remainder $\bigO(n^{-1})$ requires the Lambert phase to one order beyond a
naive term count.  Keeping $\lambda_n$ exact inside $\Psi$ avoids this bookkeeping hazard.
\end{warning}

\subsection{Compact versus expanded form}

Define
\begin{equation}
 \mathcal M_n=-\frac L2\lambda_n^2+
 \left(1+\frac L2\right)\lambda_n-\frac12\log\lambda_n
 +C_F+\Psi(\lambda_n).                                     \label{eq:Mn}
\end{equation}
Substitution of \eqref{eq:lambda-expansion} and collection of terms gives
\begin{align}
 \mathcal M_n={}&-\frac L2n^2-n\log n+
 \left(1+\frac L2\right)n-\frac{\log^2n}{2L}+C_F           \notag\\
 &-\frac{\log^2n}{2L^2n}+\Psi(\lambda_n)+\bigO(n^{-1}).   \label{eq:compact-to-expanded}
\end{align}
Thus \eqref{eq:headline-compact} and \eqref{eq:headline-expanded} are equivalent at the stated
precision.  The compact formula is structurally cleaner and is the one that survives
unchanged for arbitrary small real arguments.

\section{Sharp route II: a coefficient saddle from the exact generating product}
\label{sec:coefficient-saddle}

The endpoint--Laplace route starts from $d_n=\E(1-X)^n$.  There is a second route beginning
directly with the recurrence coefficient $a_n=[z^n]A(z)$.  It is particularly well suited to
the question posed here because $a_n$ is exactly what remains after the $q$-binomial formula
is decompressed.

\subsection{Positive-axis growth of \texorpdfstring{$A$}{A}}

The reflection identity \eqref{eq:A-reflection} and
\Cref{thm:periodic-decomposition} give, for $r>0$,
\begin{align}
 K(r):=\log A(r)
 ={}&r+\Lam(r)                                               \notag\\
 ={}&r-\frac{\log^2r}{2L}+\frac12\log r
 +P(\log_2r)+\mathcal E(r).                                \label{eq:K-positive}
\end{align}
The reference saddle is chosen by ignoring the bounded derivative of the periodic term:
\begin{equation}
 r-\frac{\log r}{L}=n.                                      \label{eq:reference-saddle}
\end{equation}
Its large solution is exactly $r=\lambda_n$.  The exact coefficient saddle for
$rK'(r)=n$ differs from $\lambda_n$ by only $\bigO(1)$, because
\begin{equation}
 rK'(r)=r-\frac{\log r}{L}+\frac12+\frac1L P'(\log_2r)
 +\bigO(r\e^{-r}).                                          \label{eq:exact-saddle-shift}
\end{equation}
An $\bigO(1)$ displacement of a saddle of curvature $\asymp1/r$ changes the logarithmic
coefficient estimate by only $\bigO(1/r)$.

\subsection{Cauchy extraction}

Cauchy's formula gives
\begin{equation}
 a_n=\frac{A(r)}{2\pi r^n}
 \int_{-\pi}^{\pi}
 \frac{A(r\e^{\ii\theta})}{A(r)}\e^{-\ii n\theta}\dd\theta.
                                                                    \label{eq:Cauchy-an}
\end{equation}
At $r=\lambda_n$, split the product \eqref{eq:A-product} near the dyadic index
$j\approx\log_2r$.  In a central window, Taylor expansion of the logarithm gives a Gaussian
with variance parameter $r+\bigO(1)$; outside that window, a positive proportion of factors
has modulus strictly below its positive-axis value, and the remaining factors supply
arbitrary polynomial decay.  The standard coefficient-saddle estimate \cite{hayman1956,flajolet2009} is therefore
\begin{equation}
 a_n=\frac{A(\lambda_n)}{\lambda_n^n\sqrt{2\pi\lambda_n}}
 \left(1+\bigO(\lambda_n^{-1})\right).                     \label{eq:an-saddle}
\end{equation}
The use of the reference rather than exact saddle is absorbed by the same relative error.

Taking logarithms and applying \eqref{eq:K-positive},
\begin{equation}
 \log a_n=\lambda_n-n\log\lambda_n
 -\frac{\log^2\lambda_n}{2L}
 +P(\log_2\lambda_n)-\frac12\log(2\pi)
 +\bigO(n^{-1}).                                            \label{eq:loga-pre}
\end{equation}
The $+\tfrac12\log\lambda_n$ in $K$ has canceled the Gaussian denominator.  Now use
$n=\lambda_n-(\log\lambda_n)/L$:
\begin{equation}
 \log a_n=\lambda_n-\lambda_n\log\lambda_n
 +\frac{\log^2\lambda_n}{2L}
 +P(\lambda_n)-\frac12\log(2\pi)
 +\bigO(n^{-1}),                                            \label{eq:loga-lambda}
\end{equation}
where periodicity was used in the form
$P(\log_2\lambda_n)=P(\lambda_n-n)=P(\lambda_n)$.

Finally, subtract $\binom n2L$ according to \eqref{eq:F-an}.  Writing
$\ell_n=\log\lambda_n$ and $n=\lambda_n-\ell_n/L$, one has
\begin{align*}
 -\binom n2L
 ={}&-\frac L2\lambda_n^2+\lambda_n\ell_n
 -\frac{\ell_n^2}{2L}+\frac L2\lambda_n-\frac12\ell_n.
\end{align*}
Adding \eqref{eq:loga-lambda} makes the terms
$\lambda_n\ell_n$ and $\ell_n^2/(2L)$ cancel exactly.  Using
\eqref{eq:CF-from-P} gives
\begin{equation}
 \log F(2^{-n})=
 -\frac L2\lambda_n^2+
 \left(1+\frac L2\right)\lambda_n-\frac12\log\lambda_n
 +C_F+\Psi(\lambda_n)+\bigO(n^{-1}).                        \label{eq:saddle-final}
\end{equation}
This is \eqref{eq:headline-compact}.

\begin{boxedremark}
The coefficient saddle provides a particularly direct answer to the original question.
Starting from the exact $q$-binomial value, \Cref{thm:q-decompression} produces $a_n$; Cauchy
extraction of that coefficient from the exact product produces the lower-Lambert asymptotic.
The endpoint--Laplace route and the coefficient route are not the same proof in disguise:
one compares $(1-X)^n$ with $\e^{-nX}$ on the real line, while the other analyzes a complex
coefficient integral.  Their agreement fixes the normalization and phase very rigidly.
\end{boxedremark}

\section{Consequences for every exact arithmetic representation}
\label{sec:corollaries}

\subsection{The normalized recurrence and composition sum}

Since $a_n=2^{\binom n2}F(2^{-n})$, \eqref{eq:headline-expanded} immediately gives
\begin{align}
 \log a_n={}&-n\log n+n-\frac{\log^2n}{2L}+C_F             \notag\\
 &-\frac{\log^2n}{2L^2n}+\Psi(\lambda_n)+\bigO(n^{-1}).   \label{eq:an-asymptotic}
\end{align}
Thus the exact composition sum \eqref{eq:composition-an} has the asymptotic
\begin{align}
 &\sum_{\rho\vDash n}
 \prod_{j=1}^{\ell(\rho)}
 \frac{1}{(2^{s_j}-1)(r_j+1)!}                             \notag\\
 &\quad=\exp\!\left(
 -n\log n+n-\frac{\log^2n}{2L}+C_F
 -\frac{\log^2n}{2L^2n}+\Psi(\lambda_n)
 +\bigO(n^{-1})\right).                                    \label{eq:composition-asymptotic}
\end{align}
The recurrence therefore has a moving periodic multiplier; no ansatz with a single fixed
constant can be sharp.

\subsection{The half moments}

Adding Stirling's expansion to \eqref{eq:an-asymptotic},
\begin{align}
 \log d_n={}&\frac12\log(2\pi n)+C_F+\Psi(\lambda_n)
 -\frac{\log^2n}{2L}                                      \notag\\
 &-\frac{\log^2n}{2L^2n}+\bigO(n^{-1}).                   \label{eq:dn-asymptotic}
\end{align}
In particular,
\begin{equation}
 d_n\sim\sqrt{2\pi n}\,
 \exp\!\left(C_F+\Psi(\lambda_n)-\frac{\log^2n}{2L}\right).
                                                                    \label{eq:dn-equivalent}
\end{equation}
The notation ``$\sim$'' is legitimate because the omitted logarithmic error tends to zero.
The periodic factor $\exp\Psi(\lambda_n)$ remains part of the equivalent.

\subsection{The literal \texorpdfstring{$q$}{q}-binomial numerator}

By \Cref{thm:q-decompression},
\begin{equation}
 \mathcal N_n=\M_na_n.                                      \label{eq:N-product}
\end{equation}
Furthermore,
\begin{equation}
 \M_n=2^{\binom{n+1}{2}}(1/2;1/2)_n                       \label{eq:M-q}
\end{equation}
and
\begin{equation}
 \log(1/2;1/2)_n=
 \log(1/2;1/2)_\infty+\bigO(2^{-n}).                       \label{eq:qpoch-limit}
\end{equation}
Consequently
\begin{align}
 \log\mathcal N_n={}&\frac L2n^2-n\log n+
 \left(1+\frac L2\right)n-\frac{\log^2n}{2L}              \notag\\
 &+C_F+\log(1/2;1/2)_\infty
 -\frac{\log^2n}{2L^2n}+\Psi(\lambda_n)
 +\bigO(n^{-1}).                                            \label{eq:N-asymptotic}
\end{align}
This formula is the sharp asymptotic of the finite signed $q$-binomial/Thue--Morse sum itself.
It makes the scale cancellation visible: the numerator grows like
$\exp(Ln^2/2)$, while the explicit denominator in \eqref{eq:q-F-N} contains
$2^{n^2}$.  The residual triangular factor is exactly the decay
$2^{-\binom n2}$ in \eqref{eq:F-an}.

\subsection{A first all-orders corollary}

The general small-argument saddle expansion supplies a first correction
\begin{equation}
 A_1(u)=\frac1{24}+\frac1{2L}
 -\frac{\Psi''(u)+\Psi'(u)^2}{2L^2}.                        \label{eq:A1-report}
\end{equation}
At reciprocal powers this gives
\begin{equation}
 \log F(2^{-n})=\mathcal M_n+
 \frac{A_1(\lambda_n)}{\lambda_n}+
 \bigO(\lambda_n^{-2}).                                    \label{eq:first-correction-dyadic}
\end{equation}
The same correction can be transported to $a_n$, $d_n$, the composition sum, and
$\mathcal N_n$ by the exact multiplicative normalizations above.  Higher orders have the same
structure: a Poincar\'e series in inverse powers of $\lambda_n$ with smooth periodic
coefficient functions evaluated at the exact phase.

\section{What each exact formula contributes asymptotically}
\label{sec:formula-audit}

It is useful to separate exact equivalence from asymptotic convenience.  The formulae compute
the same numbers, but they expose different structures.

\begingroup\small
\begin{longtable}{@{}
  >{\raggedright\arraybackslash}p{0.25\textwidth}
  >{\raggedright\arraybackslash}p{0.31\textwidth}
  >{\raggedright\arraybackslash}p{0.36\textwidth}@{}}
\toprule
Exact representation & What it exposes & Best asymptotic use\\
\midrule
\endhead
Endpoint identity
$F(2^{-n})=d_n/(n!2^{\binom n2})$ &
The triangular dyadic factor, factorial drift, and one unknown half moment &
Immediately proves the coefficient $-L/2$ of $n^2$; becomes sharp after $d_n$ is compared
with $A(-n)$\\
\addlinespace
Recurrence for $a_n$ &
A positive triangular coefficient recursion &
Reassemble into $A(2z)=E(z)A(z)$; direct iteration term by term obscures the moving saddle\\
\addlinespace
Composition sum &
A positive sum over all ordered descents of the recurrence &
Identifies $a_n$ exactly; asymptotics are best obtained after resummation into $A$, not by
selecting a few compositions\\
\addlinespace
Entire product $A(z)=\prod E(z/2^j)$ &
The dyadic scale at every analytic level &
Primary source for both the negative-Laplace harmonic sum and the coefficient saddle\\
\addlinespace
Thue--Morse block &
A zero of exact order $k$ and a finite product over dyadic scales &
Converts the inner signed power sum into a coefficient of $A(-2^kX)/A(-X)$\\
\addlinespace
Half-base $q$-binomial row &
A divided difference on $-1,-2,\ldots,-2^n$ &
Extracts the leading coefficient $a_n$ after the block cancellation; termwise asymptotics
are badly conditioned\\
\addlinespace
$q$-Pochhammer normalization &
The Mersenne product $\M_n$ in normalized form &
Contributes the elementary limit
$(1/2;1/2)_n\to(1/2;1/2)_\infty$ and cancels $\M_n$ exactly\\
\addlinespace
Denominator and valuation formulae &
Prime-adic arithmetic of the exact rationals &
Strong arithmetic information, but by themselves insufficient to determine Archimedean
magnitude or the periodic saddle phase\\
\bottomrule
\end{longtable}
\endgroup

\subsection{Why direct termwise analysis of the \texorpdfstring{$q$}{q}-sum is misleading}

The formula
\begin{equation*}
 F(2^{-n})=\frac{\mathcal N_n}{2^{n^2}(1/2;1/2)_n}
\end{equation*}
contains an explicit logarithmic contribution $-Ln^2$, whereas the final leading term is only
$-Ln^2/2$.  The missing $+Ln^2/2$ is generated inside the alternating finite numerator.
There is no contradiction: the outer Gaussian-binomial row is designed to annihilate all
lower scale powers exactly, and the surviving leading coefficient is exponentially larger
than a naive bounded-summand estimate would suggest.

The cancellation can be quantified without inspecting any individual summand.  From
\eqref{eq:N-asymptotic},
\begin{equation}
 \log\mathcal N_n=\frac L2n^2+\bigO(n\log n),               \label{eq:N-leading}
\end{equation}
while
\begin{equation}
 -\log\bigl(2^{n^2}(1/2;1/2)_n\bigr)
 =-Ln^2-\log(1/2;1/2)_\infty+\smallo(1).                   \label{eq:q-den-leading}
\end{equation}
Their sum is $-Ln^2/2+\bigO(n\log n)$, as required.

\begin{remark}[Numerical conditioning]
A floating-point implementation of the literal signed formula is normally inferior to the
positive recurrence.  It must resolve the two exact annihilations before dividing by a factor
of size $2^{n^2}$.  The decompressed identity $\mathcal N_n=\M_na_n$ is therefore not only a
proof device; it is the stable evaluation strategy.
\end{remark}

\section{How much of the full small-argument asymptotic is encoded by inverse dyadics?}
\label{sec:general-x}

\subsection{The arbitrary-argument formula}

For $x\downarrow0$, let
\begin{equation}
 \lambda(x)=-\frac1L W_{-1}(-Lx),
 \qquad
 \lambda(x)2^{-\lambda(x)}=x.                              \label{eq:lambda-x}
\end{equation}
The full endpoint saddle has the same compact main term:
\begin{equation}
 \log F(x)=
 -\frac L2\lambda(x)^2+
 \left(1+\frac L2\right)\lambda(x)-\frac12\log\lambda(x)
 +C_F+\Psi(\lambda(x))+\bigO(\lambda(x)^{-1}).              \label{eq:general-x-main}
\end{equation}
At $x=2^{-n}$ this becomes \eqref{eq:headline-compact}.  Thus the exact dyadic arithmetic
recovers the restriction of the entire arbitrary-argument asymptotic, including the same
constant and the same periodic function at the same exact phase.

\subsection{Dense sampling of the periodic correction}

The sequence of inverse dyadics does not freeze the periodic phase.  Put
\begin{equation}
 y_n=\lambda_n-n=\log_2\lambda_n.                           \label{eq:yn}
\end{equation}
Then $y_n$ is increasing and unbounded, while
\begin{equation}
 y_{n+1}-y_n\longrightarrow0.                               \label{eq:yn-increments}
\end{equation}
Indeed, $\lambda_n\sim n$ and
$y_{n+1}-y_n=\log_2(\lambda_{n+1}/\lambda_n)=\bigO(n^{-1})$.

\begin{lemma}[Density of the dyadic Lambert phases]
\label{lem:phase-density}
The fractional parts $\{\lambda_n\}=\{y_n\}$ are dense in $[0,1]$.
\end{lemma}

\begin{proof}
Fix $\theta\in[0,1]$.  For every sufficiently large integer $m$, let $n(m)$ be the first
index for which $y_{n(m)}\ge m+\theta$.  Then
\begin{equation*}
 0\le y_{n(m)}-(m+\theta)
 <y_{n(m)}-y_{n(m)-1}.
\end{equation*}
The right side tends to zero by \eqref{eq:yn-increments}.  Hence the fractional parts of
$y_{n(m)}$ tend to $\theta$.
\end{proof}

After removing the nonperiodic terms in \eqref{eq:headline-expanded}, the exact rational
sequence therefore determines $\Psi$ uniquely.  More precisely, if
\begin{align}
 R_n={}&\log F(2^{-n})+\frac L2n^2+n\log n-
 \left(1+\frac L2\right)n+\frac{\log^2n}{2L}               \notag\\
 &-C_F+\frac{\log^2n}{2L^2n},                              \label{eq:renormalized-Rn}
\end{align}
then
\begin{equation}
 R_n=\Psi(\lambda_n)+\bigO(n^{-1}).                         \label{eq:Rn-Psi}
\end{equation}
For any $u\in[0,1]$, choose a subsequence with
$\{\lambda_{n_j}\}\to u$; continuity gives $R_{n_j}\to\Psi(u)$.  In this precise sense the
inverse-power values encode the whole periodic correction, not just a single phase sample.

\subsection{What inverse powers alone do not prove}

The density result does not by itself extend \eqref{eq:headline-compact} uniformly to all
$x\downarrow0$.  Monotonicity only brackets a point
$2^{-(n+1)}\le x\le2^{-n}$ between two endpoint values whose logarithms differ by order $n$:
\begin{equation}
 \log F(2^{-n})-\log F(2^{-(n+1)})=Ln+\bigO(\log n).        \label{eq:dyadic-gap}
\end{equation}
That uncertainty is enormous compared with the desired $\bigO(n^{-1})$ remainder.

The global formula \eqref{eq:q-global} contains every dyadic rational and could, in principle,
be used on meshes much finer than one reciprocal-power interval.  To derive
\eqref{eq:general-x-main} by that route one would need a two-parameter asymptotic uniform in
both numerator and denominator.  The arbitrary-argument Bromwich saddle avoids that difficult
finite-sum uniformity.  Accordingly, the exact inverse-power formulae recover the full phase
function and all constants, but a genuinely uniform analytic argument remains necessary to
pass from the dense arithmetic set to a sharp theorem on a continuum.

\section{Numerical verification from exact rational values}
\label{sec:numerics}

The recurrence \eqref{eq:an-recurrence} computes $a_n$, hence $F(2^{-n})$, in exact rational
arithmetic.  For example,
\begin{equation}
 F(2^{-10})=
 \frac{134926369}{1246394851358539387238350848000},         \label{eq:F10}
\end{equation}
so
\begin{equation*}
 \log F(2^{-10})=-50.5775682823421608125\ldots.
\end{equation*}

Let
\begin{equation}
 \rho_n=\log F(2^{-n})-\mathcal M_n.                        \label{eq:rho}
\end{equation}
The leading theorem predicts $\rho_n=\bigO(1/\lambda_n)$, and the first correction predicts
$\lambda_n\rho_n=A_1(\lambda_n)+\bigO(1/\lambda_n)$.  The table was computed by exact rational
recurrence followed by one high-precision logarithm; $P$ and its derivatives were evaluated
from the absolutely convergent product and forward-tail series.

\begin{table}[htbp]
\centering
\small
\begin{tabular}{@{}rrrrrr@{}}
\toprule
$n$ & $\lambda_n$ & $\rho_n$ & $\lambda_n\rho_n$
& $A_1(\lambda_n)$ & $\lambda_n^2(\rho_n-A_1/\lambda_n)$\\
\midrule
10  & 13.78503056  & 0.05801333439 & 0.7997155875 & 0.7629678468 & 0.5065687288\\
20  & 24.62186833  & 0.03182756966 & 0.7836542295 & 0.7629268731 & 0.5103462408\\
40  & 45.50804986  & 0.01701388954 & 0.7742689337 & 0.7629490545 & 0.5151456258\\
80  & 86.43351899  & 0.008896709914& 0.7689739453 & 0.7629823191 & 0.5178773369\\
160 & 167.3870441   & 0.004576872612& 0.7661091776 & 0.7630070725 & 0.5192522062\\
\bottomrule
\end{tabular}
\caption{Exact reciprocal-power values against the compact Lambert main and first correction.}
\label{tab:numerics}
\end{table}

The fourth and fifth columns approach one another, while the last column remains bounded and
appears to settle, exactly as \eqref{eq:first-correction-dyadic} predicts.

A second check displays the cancellation in the literal $q$-formula.  Since
\begin{equation*}
 \log F(2^{-n})=\log\mathcal N_n-n^2L-\log(1/2;1/2)_n,
\end{equation*}
one obtains:

\begin{table}[htbp]
\centering
\small
\begin{tabular}{@{}rrrrr@{}}
\toprule
$n$ & $\log\mathcal N_n$ & $n^2L$ & $\log(1/2;1/2)_n$ & $\log F(2^{-n})$\\
\midrule
10 & 17.4960644003 & 69.3147180560 & -1.2410853733 & -50.5775682823\\
20 & 95.4684669707 & 277.2588722240 & -1.2420611411 & -180.5483441121\\
40 & 447.4055728988& 1109.0354888959& -1.2420620948 & -660.3878539023\\
80 & 1957.8744511397&4436.1419555836& -1.2420620948 & -2477.0254423491\\
\bottomrule
\end{tabular}
\caption{Logarithmic cancellation in the finite $q$-binomial formula.}
\label{tab:q-cancellation}
\end{table}

The Pochhammer logarithm has already converged to
\begin{equation}
 \log(1/2;1/2)_\infty=-1.2420620948124149458\ldots.         \label{eq:qpoch-numeric}
\end{equation}
The nontrivial compensation is between $\log\mathcal N_n\sim Ln^2/2$ and the explicit
$-Ln^2$ denominator.

\section{Correspondence with the Lean development}
\label{sec:lean-map}

The relevant modules in the 25 August 2026 repository snapshot \cite{proveit2026}
form a fairly literal proof graph.  The table records the role of each group rather than every import edge.

\begingroup\small
\begin{longtable}{@{}
  >{\raggedright\arraybackslash}p{0.33\textwidth}
  >{\raggedright\arraybackslash}p{0.59\textwidth}@{}}
\toprule
Module & Role in the bridge\\
\midrule
\endhead
\path{FabiusRecurrenceSequence.lean} &
Defines $a_n=d_n/n!$, proves its recurrence, identifies
$F(2^{-n})=2^{-\binom n2}a_n$, and connects its ordinary series to the analytic generating
function and dyadic product.\\
\addlinespace
\path{FabiusInverseDyadicClosedForm.lean} &
Unfolds the inverse-dyadic recurrence into the exact composition formula.\\
\addlinespace
\path{FabiusRawQBinomialFormula.lean},
\path{FabiusQBinomialFormula.lean},
\path{FabiusQBinomialScalarFormula.lean},
\path{FabiusDyadicQBinomialScalar.lean} &
Develop the Thue--Morse block factorization, half-base Gaussian interpolation, scalar
translation invariance, and the global exact dyadic formula.  The coefficient-extraction
argument in \Cref{sec:q-decompression} is the ordinary-mathematics unpacking of this layer.\\
\addlinespace
\path{FabiusDyadicLogBounds.lean} &
Combines the exact endpoint identity with elementary half-moment bounds to prove the
$-Ln^2/2$ law and an explicit $3n\log(n+1)$ error.\\
\addlinespace
\path{NegativeLaplace.lean},
\path{PeriodicCorrection.lean},
\path{PeriodicFourier.lean} &
Construct the negative-Laplace logarithm, its exact quadratic-plus-periodic decomposition,
the forward-tail bound, regularity, mean normalization, and Fourier data.\\
\addlinespace
\path{EndpointLaplaceComparison.lean} &
Proves the second-order pointwise and integrated comparison between $(1-x)^n$ and
$\e^{-nx}(1-nx^2/2)$ and transfers the estimate through the logarithm.\\
\addlinespace
\path{DyadicSharpDecomposition.lean},
\path{FabiusDyadicSharpCumulant.lean} &
Assemble the exact logarithmic decomposition of $F(2^{-n})$, including Stirling, the periodic
term, the endpoint/Laplace error, and the exact second tilted cumulant.\\
\addlinespace
\path{LaplacePeriodicSecondOrder.lean} &
Differentiates the periodic Laplace decomposition and proves the explicit
$\Psi'(\log_2n)$ correction with an $\bigO(1/n)$ remainder.\\
\addlinespace
\path{FabiusLambertPhase.lean},
\path{FabiusDyadicSharpAsymptotic.lean} &
Construct the exact lower-Lambert phase, transfer the periodic derivative term to
$\Psi(\lambda_n)$, and prove the corrected sharp dyadic theorem.\\
\addlinespace
\path{FabiusSharpLambertMain.lean},
\path{FabiusSharpAsymptotic.lean},
\path{FabiusFullAsymptoticExpansion.lean} &
Prove the arbitrary-small-argument Lambert main and continue the saddle expansion to all
orders.\\
\bottomrule
\end{longtable}
\endgroup

The coefficient-saddle derivation in \Cref{sec:coefficient-saddle} is deliberately marked as
an alternative exposition.  The formal dyadic proof uses the endpoint--Laplace comparison;
the formal arbitrary-argument proof uses the Bromwich/saddle kernel.  No claim is made that
the Cauchy coefficient route appears as a separate theorem under that exact name.

\section{Conclusions}
\label{sec:conclusion}

The conspicuous gap between exact dyadic evaluation and endpoint asymptotics can be closed
without adding a new special function or an independent heuristic ansatz.

First, the exact reciprocal-power values are controlled by a single coefficient sequence:
\begin{equation*}
 F(2^{-n})=2^{-\binom n2}a_n,
 \qquad
 A(z)=\sum_{n\ge0}a_nz^n=\prod_{j\ge1}E(z/2^j).
\end{equation*}
The recurrence and composition formula are simply two expansions of this identity.

Second, the $q$-binomial formula reduces exactly to the same sequence.  The inner Thue--Morse
sum is a coefficient of $A(-2^kX)/A(-X)$, and the outer half-base Gaussian row is the
leading-coefficient functional on the nodes $-2^k$.  Therefore
\begin{equation*}
 \mathcal N_n=\M_na_n
\end{equation*}
with no approximation.  This is the direct algebraic connection that makes asymptotic
analysis of the finite $q$-formula feasible.

Third, the negative-axis product of $A$ satisfies an exact dyadic dilation equation.  Removing
its quadratic drift leaves a smooth periodic function.  A second-order comparison between the
endpoint kernel $(1-X)^n$ and the Laplace kernel $\e^{-nX}$ supplies the first cumulant
correction; differentiating the exact periodic decomposition identifies that correction; and
the lower-Lambert fixed point moves it to the exact phase.  The result is
\begin{equation*}
 \log F(2^{-n})=
 -\frac L2\lambda_n^2+
 \left(1+\frac L2\right)\lambda_n-\frac12\log\lambda_n
 +C_F+\Psi(\lambda_n)+\bigO(n^{-1}).
\end{equation*}

Finally, the inverse dyadic sequence is richer than a sparse one-phase sample: the fractional
parts of $\lambda_n$ are dense, so the exact rational values determine the entire continuous
periodic correction after renormalization.  What they do not provide by monotonicity alone is
a uniform continuum theorem at the sharp error scale; that step still requires a uniform
saddle argument or a two-parameter analysis of the general dyadic formula.

The answer to the motivating question is therefore affirmative in a strong sense.  The
small-argument asymptotic can be re-derived from the exact calculation formulae, including the
constant, the logarithmic correction, the exact lower-Lambert phase, and the nonconstant
periodic fluctuation.  The only essential move is to respect the exact cancellations built
into those formulae rather than attempting to estimate their finite summands independently.

\appendix

\section{Algebra behind the expanded dyadic formula}
\label{app:algebra}

This appendix records the expansion needed to pass from the compact Lambert form to the
explicit expression in $n$.  Put $\ell=\log n$ and seek
\begin{equation}
 \lambda_n=n+\frac\ell L+\frac{c_1}{n}+\frac{c_2}{n^2}
 +\bigO(\ell^3/n^3).                                       \label{eq:lambda-ansatz}
\end{equation}
The fixed-point equation is
\begin{equation}
 \lambda_n=n+\frac1L\log\lambda_n.                         \label{eq:lambda-fixed-app}
\end{equation}
Writing
\begin{equation*}
 \lambda_n=n\left(1+\frac\ell{Ln}+\frac{c_1}{n^2}
 +\frac{c_2}{n^3}+\cdots\right)
\end{equation*}
and expanding the logarithm gives
\begin{equation*}
 \log\lambda_n=
 \ell+\frac\ell{Ln}+\frac{c_1-\ell^2/(2L^2)}{n^2}
 +\bigO(\ell^3/n^3).
\end{equation*}
Comparison in \eqref{eq:lambda-fixed-app} yields
\begin{equation}
 c_1=\frac\ell{L^2},
 \qquad
 c_2=\frac{\ell-\ell^2/2}{L^3},                            \label{eq:c1-c2}
\end{equation}
which gives the first four terms of \eqref{eq:lambda-expansion}.

Now set
\begin{equation*}
 H(\lambda)=-\frac L2\lambda^2+
 \left(1+\frac L2\right)\lambda-\frac12\log\lambda.
\end{equation*}
Substitution of \eqref{eq:lambda-ansatz} with \eqref{eq:c1-c2} and collection through order
$1/n$ gives
\begin{align}
 H(\lambda_n)={}&-\frac L2n^2-n\ell+
 \left(1+\frac L2\right)n-\frac{\ell^2}{2L}                \notag\\
 &-\frac{\ell^2}{2L^2n}+\bigO(n^{-1}).                    \label{eq:H-expanded}
\end{align}
The notation $\bigO(n^{-1})$ in \eqref{eq:H-expanded} includes bounded multiples of $1/n$;
the displayed $\ell^2/n$ term is separated because it is asymptotically larger.  Adding
$C_F+\Psi(\lambda_n)$ proves \eqref{eq:compact-to-expanded}.

\section{A stable numerical recipe}
\label{app:numerical}

The following procedure reproduces \Cref{tab:numerics} without evaluating the unstable signed
$q$-sum.

\begin{algorithm}[Exact arithmetic plus convergent products]
\label{alg:numerics}
\leavevmode
\begin{enumerate}
\item Compute $a_0=1$ and, for $1\le j\le n$,
\begin{equation*}
 a_j=\frac1{2^j-1}\sum_{k<j}\frac{a_k}{(j-k+1)!}
\end{equation*}
in rational arithmetic.

\item Form the exact rational
\begin{equation*}
 F(2^{-n})=2^{-\binom n2}a_n
\end{equation*}
and take its logarithm only after conversion to the desired floating-point precision.

\item Solve $\lambda-\log\lambda/L=n$ on the large branch, either by
$\lambda=-W_{-1}(-L2^{-n})/L$ or by Newton iteration initialized with
$n+\log n/L$.

\item Reduce $u=\lambda-\lfloor\lambda\rfloor$ and put $s=2^u\in[1,2)$.  Evaluate
\begin{align*}
 P(u)={}&\sum_{j\ge1}\log\frac{1-\e^{-s/2^j}}{s/2^j}
 +\frac L2(u^2-u)\\
 &+\sum_{m\ge0}\log(1-\e^{-s2^m}).
\end{align*}
The first series has a geometric small-argument tail and the second a doubly exponential tail.
Use stable primitives such as \texttt{expm1} for the first factors.

\item Subtract $\overline P$ from \eqref{eq:P-mean} to obtain $\Psi(u)$.  Differentiate the
locally uniformly convergent series termwise, or use the rapidly convergent Fourier series
\eqref{eq:Psi-Fourier-report}, to obtain $\Psi'$ and $\Psi''$.

\item Evaluate $\mathcal M_n$ from \eqref{eq:Mn} and, when desired, $A_1$ from
\eqref{eq:A1-report}.
\end{enumerate}
\end{algorithm}

For a direct verification of the $q$-formula one may compute
$\mathcal N_n=\M_na_n$ exactly and compare it with the finite signed expression.  The equality
is a better test than comparing floating-point approximations of the two sides because it
checks both cancellation layers over the rationals.

\section{Notation cross-reference}
\label{app:notation}

\begingroup\small
\begin{longtable}{@{}ll@{}}
\toprule
Symbol & Meaning\\
\midrule
\endhead
$L$ & $\log2$\\
$F$ & bounded Fabius function / distribution function on $[0,1]$\\
$\Up$ & Rvachev up-function, with $\Up(t)=F(1-|t|)$ on $[-1,1]$\\
$X$ & random variable with distribution function $F$\\
$d_n$ & half moment $\E(1-X)^n$\\
$a_n$ & factorially normalized half moment $d_n/n!$\\
$A(z)$ & $\sum a_nz^n=\E\e^{zX}$\\
$E(z)$ & $(\e^z-1)/z$\\
$\M_n$ & $\prod_{j=1}^n(2^j-1)$\\
$\mathcal N_n$ & finite $q$-binomial/Thue--Morse numerator for $F(2^{-n})$\\
$B(s)$ & $A(-s)=\E\e^{-sX}$\\
$\Lam(s)$ & $\log B(s)$\\
$P$ & exact one-periodic correction in $\Lam$\\
$\Psi$ & $P-\overline P$, mean-zero periodic correction\\
$C_F$ & constant $\overline P-\tfrac12\log(2\pi)$\\
$\lambda_n$ & large solution of $\lambda-\log\lambda/L=n$\\
$\mathcal M_n$ & compact Lambert main term for $\log F(2^{-n})$\\
\bottomrule
\end{longtable}
\endgroup

\begin{thebibliography}{99}

\bibitem{arias1982}
J.~Arias de Reyna,
\newblock Definici\'on y estudio de una funci\'on indefinidamente diferenciable de soporte
  compacto,
\newblock \emph{Revista de la Real Academia de Ciencias Exactas, F\'isicas y Naturales de
  Madrid} \textbf{76} (1982), no.~1, 21--38;
  English translation, arXiv:\nolinkurl{1702.05442} (2017).

\bibitem{arias2018}
J.~Arias de Reyna,
\newblock Arithmetic of the Fabius function,
\newblock \emph{INTEGERS} \textbf{18} (2018), Article A51;
  arXiv:\nolinkurl{1702.06487}, version~3.

\bibitem{andrews1976}
G.~E.~Andrews,
\newblock \emph{The Theory of Partitions},
\newblock Encyclopedia of Mathematics and its Applications~2,
  Addison--Wesley, 1976.

\bibitem{allouche2003}
J.-P.~Allouche and J.~Shallit,
\newblock \emph{Automatic Sequences: Theory, Applications, Generalizations},
\newblock Cambridge University Press, 2003.

\bibitem{corless1996}
R.~M.~Corless, G.~H.~Gonnet, D.~E.~G.~Hare, D.~J.~Jeffrey, and D.~E.~Knuth,
\newblock On the Lambert $W$ function,
\newblock \emph{Advances in Computational Mathematics} \textbf{5} (1996), 329--359.

\bibitem{dlmf17}
NIST Digital Library of Mathematical Functions, Chapter~17,
\newblock \emph{$q$-Hypergeometric and Related Functions},
\newblock \url{https://dlmf.nist.gov/17}.

\bibitem{flajolet1995}
P.~Flajolet, X.~Gourdon, and P.~Dumas,
\newblock Mellin transforms and asymptotics: harmonic sums,
\newblock \emph{Theoretical Computer Science} \textbf{144} (1995), 3--58.

\bibitem{flajolet2009}
P.~Flajolet and R.~Sedgewick,
\newblock \emph{Analytic Combinatorics},
\newblock Cambridge University Press, 2009.

\bibitem{gasper2004}
G.~Gasper and M.~Rahman,
\newblock \emph{Basic Hypergeometric Series}, second edition,
\newblock Encyclopedia of Mathematics and its Applications~96,
  Cambridge University Press, 2004.

\bibitem{hayman1956}
W.~K.~Hayman,
\newblock A generalisation of Stirling's formula,
\newblock \emph{Journal f\"ur die reine und angewandte Mathematik}
  \textbf{196} (1956), 67--95.

\bibitem{haugland2016}
J.~K.~Haugland,
\newblock Evaluating the Fabius function,
\newblock arXiv:\nolinkurl{1609.07999} (2016).

\bibitem{kac2002}
V.~Kac and P.~Cheung,
\newblock \emph{Quantum Calculus},
\newblock Universitext, Springer, 2002.

\bibitem{reshetnikovMSE}
V.~Reshetnikov,
\newblock A conjectured closed form for the Fabius function at dyadic rationals,
\newblock Mathematics Stack Exchange, question 3283519, 5 July 2019.

\bibitem{proveit2026}
V.~Reshetnikov,
\newblock \emph{ProveIt: Analysis/FabiusFunction},
\newblock Lean~4 formal development, repository snapshot
  \path{a24af8347aba5d38b8febf2cf9a19eeef6aba18a}, 25 August 2026.
\newblock
\url{https://github.com/VladimirReshetnikov/ProveIt/tree/a24af8347aba5d38b8febf2cf9a19eeef6aba18a/Analysis/FabiusFunction}.

\bibitem{fabiusExposition2026}
V.~Reshetnikov,
\newblock \emph{The Fabius Function and Rvachev's Up-Function},
\newblock technical exposition in the same repository snapshot, 25 August 2026.

\end{thebibliography}

\end{document}
